\section{Related work}
\label{sec:related}

\paragraph{Quantum machine learning for financial prediction.}
Reports of quantum models for forecasting and stock selection range from
clear positives on synthetic data --- a 96-parameter circuit matching
thousands-parameter BiLSTMs on noisy sinusoids \citep{emmanoulopoulos2022quantum}
--- to mixed results on market data: a hybrid network beats a parameter-matched
ANN on directional classification of single stocks, yet its recurrent variant
wins on risk-adjusted returns in only two of four S\&P~500 regimes under
identical walk-forward splits \citep{ahmad2026quantum}; on OHLCV data no
model beats the naive baseline intraday and quantum models merely match
classical ones daily \citep{shah2026stock}; and adding quantum layers to a
Transformer degrades performance on most tasks \citep{chen2025quantumd}. The
two studies closest to ours are both controlled negatives on the equity
cross-section. On TOPIX500 over a ten-year rolling backtest the quantum neural
network is inferior to its classical counterpart over the whole period
\citep{kobayashi2023cross}. On China A-shares, a kernel-swap design that holds
subsamples, solver and tuning budget fixed finds the fidelity quantum kernel
indistinguishable from an RBF control ($\Delta\mathrm{IC}=+0.005$, $p=0.42$)
over 170 walk-forward windows, with no pairwise difference among eleven models
surviving family-wise correction --- and shows that a 60-window evaluation on a
hindsight-screened universe makes the same kernel look dominant
\citep{shen2026quantum}. That paper's diagnosis of the applied literature ---
short samples, hindsight universes, weak baselines --- and its proposed
standards (kernel-swap controls, budget-equalised comparisons, point-in-time
universes, multiplicity-robust inference) are the ones our protocol adopts. The
reference QTCNN study we audit \citep{chen2025quantumb} reports an annualised
out-of-sample Sharpe ratio of \num{0.538} $\pm$ \num{0.042} over a 22-day
window, acknowledges that 22 days give limited power, and defers a full-sample
backtest to future work; its own encoder-only ablation scores \num{0.412},
which already places most of the reported edge in the classical encoder.

\paragraph{QCNN trainability and dequantisation.}
Randomly initialised parametrised circuits generically exhibit barren
plateaus \citep{mcclean2018barren,larocca2024barren}; global cost functions
induce them even at shallow depth whereas local observables keep gradients
polynomial up to depth $O(\log n)$ \citep{cerezo2021cost}, hardware noise
induces its own plateau \citep{wang2021noise}, and data embeddings add a
dataset-induced variant \citep{thanasilp2021subtleties}. The QCNN of
\citet{cong2019quantum} avoids the problem: its gradient variance vanishes at
most polynomially, and sharing parameters across a layer, as in the original
design, enlarges the variance further \citep{pesah2020absence}. Our
gradient-variance and expressibility scans (Section~\ref{sec:results}) are the
empirical counterpart of these guarantees for the ansatz we use. The guarantee
cuts both ways. Randomly initialised QCNNs act only on low-weight information
about their inputs, and a classical surrogate built from Pauli propagation
matches or beats them on every benchmark up to 1024 qubits
\citep{bermejo2024quantum}; more generally, the structure that removes barren
plateaus tends to confine the model to a classically simulable subspace
\citep{cerezo2023provable}. Classical surrogates that reproduce a trained
circuit's input--output map are therefore the natural benchmark a quantum model
must beat \citep{schreiber2022classical}; random-Fourier-feature surrogates can
be built from the architecture alone \citep{landman2022classically}; wide
random circuits converge to Gaussian processes and stay there under training
\citep{girardi2024trained}; and in the kernel view the optimal model is a convex
problem \citep{schuld2021quantum}. Our \textsc{Frozen} and \textsc{Surrogate}
conditions operationalise this random-feature reading: if the trained circuit
is no better than a frozen one, and the frozen one no better than a classical
map of equal parameter count, the layer is a random feature and the
dequantisation argument applies without further analysis. Related diagnostics
warn against reading performance as evidence of quantum effects: hybrid
networks show no consistent expressibility--trainability relation
\citep{kashif2026rethinking}, a fixed encoding can impose a capacity ceiling
that mimics a plateau \citep{pandey2026architectural}, current models fit random
labels so uniform generalisation bounds are vacuous
\citep{gilfuster2023understanding}, and a causal-mediation analysis finds
quantum-mediated effects averaging \SI{0.82}{\percent} with \SI{93}{\percent} of
configurations in a neutral regime \citep{kesiku2026quantum}.

\paragraph{Benchmarking practice in hybrid quantum ML.}
Systematic benchmarks find prototypical classical baselines beating
prototypical quantum models, with entanglement removable at no cost and only
one paper in the authors' literature sample drawing critical conclusions from
its own results \citep{bowles2024better}; an industrial comparison finds
classical and hybrid models statistically equivalent, with accuracy
uncorrelated with entangling power \citep{basilewitsch2024quantum}. Budget-aware
protocols that match trainable parameters and optimisation steps find shallow
hybrids matching classical performance at much higher cost
\citep{bendouma2026fair}, whereas a same-accuracy-threshold comparison reports
hybrids scaling FLOPs more efficiently \citep{kashif2024computational}. The
attribution audit of \citet{gillani2026quantum} --- parameter-matched classical
controls sharing the identical front-end, a random-feature kernel, a
regularisation sweep, three to five seeds with confidence intervals, and the
explicit stance that the benchmark's value does not depend on quantum winning
--- is the closest methodological template to ours; the capacity-matched
bottleneck of \citet{chen2025quantumd} follows the same commitment. Shot
budgets are the operative hardware cost, since access is priced per shot
\citep{liu2025you} and gradient estimation needs two circuit runs per parameter
per iteration \citep{scriva2023challenges,miroszewski2024search}; batched
parameter broadcasting is what removes per-call simulator overhead
\citep{chen2025batched}.

\paragraph{Evaluation methodology in empirical finance.}
A handful of trials suffices to find a spuriously high backtested Sharpe ratio
\citep{bailey2014pseudo}, calibrating rules by backtest itself overfits
\citep{carr2014determining}, and standard holdout is unreliable, motivating the
Deflated Sharpe Ratio \citep{bailey2014deflated} and the combinatorially
symmetric estimate of the probability of backtest overfitting
\citep{bailey2017probability} that we report. Adaptive specification search
produces significant backtests under martingale-difference nulls, and effective
multiplicity rather than the nominal specification count sets the expected
in-sample maximum \citep{nikolopoulos2026spurious}; publication selects for
author and sampling error alike \citep{chen2022publication}. Multiple-testing
control in the cross-section is contested in magnitude --- 93\% of
trading-friction anomalies fail replication and microcaps drive equal-weighted
results in \citet{hou2020replicating}, while a Bayesian replication finds the
majority of factors survive \citep{jensen2021there} --- but not in principle,
which is why we apply Romano--Wolf step-down control \citep{romano2005stepwise}
across every contrast. Leakage inflates out-of-sample performance while
preserving the appearance of rigour \citep{bhand2026illusion}; a one-switch
benchmark shows chronological splitting alone is insufficient and fixes
decision-time semantics as features known by the close of day~$t$ and execution
at the $t{+}1$ open \citep{zhang2026alpha}, which is the convention our
target definition and embargo follow. Transaction costs are decisive at daily
turnover: technical rules that pass an SPA test in sample do not survive
moderate costs \citep{rink2023predictive}, 73.3\% directional accuracy
coexists with a post-cost loss at 471 trades per year \citep{pall2026transaction},
and a systematic falsification of fourteen intraday signal families finds none
passing a $t\ge 2$, $\ge 30$-trade, net-positive standard
\citep{mesfin2026structural}. Our cost curves and break-even analysis are the
corresponding checks for a cross-sectional long--short book.
